\section{Effective compactness and local-to-global constructions}
\label{sec:compact}

\subsection{Finite nets in place of abstract compactness}

Let \((X,d)\) be represented at finite precision.  An \(\varepsilon\)-net is a
finite list \(x_0,\ldots,x_{N-1}\) together with a coverage certificate
\[
       \forall x\in X\;\exists j<N\quad d(x,x_j)\le\varepsilon.
\]
In iPRA the displayed quantifiers are not taken over an abstract completed
space.  The presentation supplies a finite observation relation strong enough
to verify the coverage statement needed by the theorem.  For a rational box,
for example, a grid and coordinate rounding algorithm provide the certificate.
For a compact metric presentation, the net producer is part of the input
interface.

\subsection{Approximate minimization}

Suppose \(f:X\to\Q\) has a modulus
\[
 d(x,y)\le\varepsilon\quad\Longrightarrow\quad
 |f(x)-f(y)|\le\omega(\varepsilon).
\]
Evaluate \(f\) on the finite net and choose the least index \(j_*\) attaining
the minimum rational value.

\begin{theorem}[Finite-net minimum transfer]
For every represented \(x\in X\),
\[
       f(x_{j_*})\le f(x)+2\omega(\varepsilon)+2e,
\]
where \(e\) bounds each finite evaluation error.  In particular, the selected
net point is an explicit approximate minimizer.
\end{theorem}

\begin{proof}
Choose a net point \(x_j\) within \(\varepsilon\) of \(x\).  Let
\(\widetilde f_j\) denote the computed rational approximation, with
\(|\widetilde f_j-f(x_j)|\le e\).  By finite minimality,
\(\widetilde f_{j_*}\le\widetilde f_j\).  Hence
\[
 f(x_{j_*})
 \le \widetilde f_{j_*}+e
 \le \widetilde f_j+e
 \le f(x_j)+2e
 \le f(x)+\omega(\varepsilon)+2e.
\]
The stated slightly looser bound allows asymmetric input/output moduli and is
closed under the usual error-budget composition.  Every choice is bounded
least search.
\end{proof}

\subsection{Finite covers}

An effective finite cover consists of:
\begin{itemize}
\item finitely many chart or open-set codes \(U_i\);
\item a Boolean membership or positive-margin predicate at the required
precision;
\item a coverage certificate for the finite observation domain;
\item transition data on pairwise overlaps;
\item a Lebesgue-number or refinement margin when exact chart selection is
needed.
\end{itemize}
The framework never extracts a finite subcover from an arbitrary open cover.
It consumes or constructs a finite cover from explicit quantitative data.

\subsection{Finite gluing}

Let local data \(s_i\) be given on a finite cover.  Suppose compatibility on
overlaps is a bounded decidable relation
\[
       R_{ij}(s_i,s_j)=1.
\]
A gluing theorem can have one of three outputs:
\begin{enumerate}
\item an exact global code, when a bounded algebraic constructor exists;
\item a finite cloud of global codes, with quotient-stable observable;
\item a regular name obtained from compatible approximate gluings and a
shrinking diameter bound.
\end{enumerate}

For affine algebraic data, gluing often reduces to a finite system of ring
equalities.  For simplicial data, it reduces to finite carrier and contiguity
relations.  For analytic local solutions, it reduces to rational coefficient
bounds and overlap estimates.

\begin{theorem}[Finite gluing by bounded search]
Assume a primitive-recursive bound \(B\) on global candidate codes and a
bounded verifier \(G(c,s_0,\ldots,s_{r-1})\) that is sound and complete for the
finite gluing problem.  Then the finite list of all global gluings is primitive
recursive.  If a stable observable is constant on that list, it is a
well-defined deterministic output.
\end{theorem}

\begin{proof}
Enumerate \(c<B\) and retain exactly those passing \(G\).  Soundness says every
retained code is a gluing; completeness says every gluing appears.  The list is
a bounded fold.  Constancy of the observable erases the choice of retained
code.  The complete list, rather than an unbounded existential, is the
constructive global object.
\end{proof}

\subsection{Local differential equations}

For a rational polynomial ODE on a rational box, a finite local solution may
be represented by a truncated power series plus a remainder bound.  A
contracted local solver consumes:
\begin{itemize}
\item coefficient bounds;
\item a Lipschitz or majorant constant;
\item a rational step size satisfying the contraction inequality;
\item an initial-value approximation and error budget.
\end{itemize}
It returns a finite coefficient table and a verified residual bound.  A global
solution on a finite time interval is obtained by a bounded number of local
steps.  The state invariant accumulates the error budget; no abstract Picard
fixed point is needed in the object language.

\subsection{Local algebraic branches}

An equation \(F(x,y)=0\) with a certified nonzero partial derivative
\(|\partial_yF|\ge\eta\) admits a guarded local branch.  A rational Newton or
bisection procedure produces a finite regular name; the derivative margin
controls uniqueness and demand propagation.  At singular points the framework
returns a finite relation of branches or consumes exact algebraic splitting
data.  This distinction addresses the recurring Puiseux difficulty: exact
branch coefficients are safe for exact algebraic presentations, but not as a
uniform operation on arbitrary Cauchy coefficients.

\subsection{Observational quotients and extensionality}

Suppose two presentation codes \(x,x'\) represent the same object by
\(x\equiv_p x'\) at all precisions.  A contracted transducer is extensional if
for every output query \(q\), input equivalence at \(\Need_F(q)\) implies
output equivalence at \(q\).  Composition preserves extensionality by
substituting the demand inequalities.  Thus quotient reasoning becomes an
ordinary theorem about finite observations; no quotient type or choice of
canonical representatives is required inside iPRA.

\subsection{Local-to-global principle}

The general local-to-global pattern is now finite:
\begin{enumerate}
\item produce a finite cover or decomposition;
\item solve each local problem with explicit contracts;
\item verify pairwise compatibility by bounded formulas;
\item enumerate or construct finite global candidates;
\item erase nonuniqueness by exact stability or a diameter modulus;
\item compile the entire bounded process with one trace invariant.
\end{enumerate}
This pattern applies to finite sheaves, charted manifolds, local analytic
branches, and many descent calculations.  What it does not do is turn an
ineffective existence theorem for a cover, bound, or local solution into an
effective producer.
